% Standalone problem document, generated from the adjacent README.md.
% Compile with XeLaTeX. Mathematical content is maintained in Markdown.
\documentclass[11pt,a4paper]{article}
\usepackage[fontsize=10.5pt]{scrextend}
\usepackage[margin=22mm,headheight=15pt,headsep=9mm,footskip=11mm]{geometry}
\usepackage{fontspec}
\setmainfont{texgyrepagella}[Extension=.otf,UprightFont=*-regular,BoldFont=*-bold,ItalicFont=*-italic,BoldItalicFont=*-bolditalic]
\setsansfont{texgyreheros}[Extension=.otf,UprightFont=*-regular,BoldFont=*-bold,ItalicFont=*-italic,BoldItalicFont=*-bolditalic]
\setmonofont{texgyrecursor}[Extension=.otf,UprightFont=*-regular,BoldFont=*-bold,ItalicFont=*-italic,BoldItalicFont=*-bolditalic,Scale=MatchLowercase]
\usepackage{amsmath,amssymb}
\usepackage{unicode-math}
\setmathfont{texgyrepagella-math.otf}
\usepackage{xcolor,microtype,enumitem,needspace,fancyhdr}
\usepackage{bookmark}
\definecolor{ink}{HTML}{17344B}
\definecolor{accent}{HTML}{197D80}
\definecolor{muted}{HTML}{596773}
\hypersetup{colorlinks=true,linkcolor=accent,urlcolor=accent,pdftitle={MI-23 - Corrected
generalized geometric-mean product
conjecture},pdfauthor={Open Problems in Numerical Linear Algebra}}
\urlstyle{same}
\setlength{\parindent}{0pt}
\setlength{\parskip}{5pt plus 1pt minus 1pt}
\linespread{1.04}
\setlength{\emergencystretch}{3em}
\widowpenalty=10000
\clubpenalty=10000
\displaywidowpenalty=10000
\allowdisplaybreaks[1]
\setlist{nosep,leftmargin=1.5em}
\providecommand{\tightlist}{\setlength{\itemsep}{0pt}\setlength{\parskip}{0pt}}
\providecommand{\pandocbounded}[1]{#1}
\pagestyle{fancy}
\fancyhf{}
\fancyhead[L]{\sffamily\footnotesize\color{muted}OPEN PROBLEMS IN NUMERICAL LINEAR ALGEBRA}
\fancyhead[R]{\sffamily\bfseries\footnotesize\color{accent}MI-23}
\fancyfoot[L]{\parbox[b]{0.9\textwidth}{\sffamily\fontsize{8}{10}\selectfont\color{muted}Editorial ratings; a bounded search cannot certify that a problem remains unsolved.\\Literature check: 2026-09-12}}
\fancyfoot[R]{\sffamily\footnotesize\color{muted}\thepage}
\renewcommand{\headrulewidth}{0.3pt}
\renewcommand{\headrule}{\hbox to\headwidth{\color{accent}\leaders\hrule height \headrulewidth\hfill}}
\makeatletter
\renewcommand{\section}{\Needspace{5\baselineskip}\@startsection{section}{1}{0pt}{12pt plus 2pt minus 2pt}{4pt}{\sffamily\bfseries\large\color{ink}}}
\renewcommand{\subsection}{\Needspace{4\baselineskip}\@startsection{subsection}{2}{0pt}{9pt plus 1pt minus 1pt}{3pt}{\sffamily\bfseries\normalsize\color{ink}}}
\makeatother
\setcounter{secnumdepth}{0}
\begin{document}
{\sffamily\small\color{muted}Matrix inequalities and norms\par}
\vspace{3pt}
{\raggedright\hyphenpenalty=10000\exhyphenpenalty=10000\sffamily\bfseries\fontsize{21}{24}\selectfont\color{ink}Corrected
generalized geometric-mean product conjecture\par}
\vspace{7pt}
{\sffamily\small\textbf{Difficulty:} challenging\quad\textbf{Importance:} interesting
to specialist\par}
{\sffamily\small\bfseries\color{accent}Status: Lean verified\par}
\vspace{4pt}
{\color{accent}\hrule height 0.8pt}
\vspace{6pt}
\textbf{Rating rationale:} The four interacting exponents make this
corrected conjecture challenging; its immediate consequences concern
specialists in matrix means and log-majorization.

\subsection{Lean proof and verification evidence -
2026-09-12}\label{lean-proof-and-verification-evidence---2026-09-12}

\textbf{The complete corrected MI-23 conjecture is false, with a
Lean-verified counterexample.} The rational complex positive-definite
witness has \(r=s=1\), \(p=2\), \(t=1/8\) and violates the first ordered
eigenvalue inequality. The
\href{https://github.com/sgstepaniants/OpenProblemsInNLA/tree/17194f9060609acae429e14d3dc3c4562b84f2bd/matrix-inequalities-and-norms/MI-23/lean}{proof
at revision 17194f9} preserves every original dimension, both real
\(r,s\) regions, actual CFC matrix powers, all proper prefix products
and equality of complete products.

\textbf{Mathematical counterexample and informal proof:} Matthew J.
Colbrook, Department of Applied Mathematics and Theoretical Physics,
University of Cambridge. \textbf{Lean formalization:} George
Stepaniants, Department of Computing and Mathematical Sciences,
California Institute of Technology, Pasadena, California, USA, with
AI-agent assistance.

The eight
\href{https://github.com/sgstepaniants/OpenProblemsInNLA/blob/17194f9060609acae429e14d3dc3c4562b84f2bd/matrix-inequalities-and-norms/MI-23/lean/Solution.lean}{checked
exports}, each with prefix \textit{NLA.MI23.}, are:

\begin{itemize}
\tightlist
\item
  \textit{positive\_powers\_and\_means}: genuine CFC powers and
  generalized means are positive definite.
\item
  \textit{product\_eigenvalue\_semantics}: complete characteristic roots
  with multiplicities, positivity, ordering, similarity and determinant
  product.
\item
  \textit{squared\_product\_largest}: the actual largest root of
  \(X^2Y^2\) equals \(\|XY\|_2^2\).
\item
  \textit{operator\_norm\_bounds}: entry and Frobenius bounds for the
  genuine Euclidean operator norm.
\item
  \textit{witness\_data}: admissibility and all actual CFC witness
  identities.
\item
  \textit{witness\_squared\_gap}: the exact rational gap and strict
  operator-norm separation.
\item
  \textit{counterexample}: failure of the original log-majorization
  relation at the admissible witness.
\item
  \textit{not\_generalizedGeometricMeanConjecture}: negation of the
  complete universal conjecture.
\end{itemize}

Two independent agents reviewed the
\href{https://github.com/ajt60gaibb/OpenProblemsInNLA/blob/main/matrix-inequalities-and-norms/MI-23/lean/reviews}{frozen
statements and completed proof}.
\href{https://github.com/sgstepaniants/OpenProblemsInNLA/actions/runs/34716784038}{Linux
run 34716784038} matched all eight exports using the sandboxed
Comparator and replayed the solution in Lean's default kernel. The
\href{https://github.com/ajt60gaibb/OpenProblemsInNLA/blob/main/matrix-inequalities-and-norms/MI-23/lean/verification/linux-2026-09-12}{original
artifacts and independent operational audit} bind all 133 verified
source inputs and both actual isolation and rejection-control suites.
All 64
\href{https://github.com/ajt60gaibb/OpenProblemsInNLA/blob/main/matrix-inequalities-and-norms/MI-23/lean/verification/linux-2026-09-12/axiom-verification.json}{internal/public
transitive axiom reports} contain only \textit{propext},
\textit{Classical.choice} and \textit{Quot.sound}. These are independent
agent reviews, without a claim of external human peer review.

The proof pins \textbf{Lean 4.33.1},
\href{https://github.com/alerad/leancert/tree/621a43d7cf21f87872392a01e874f2f1dbddc926}{LeanCert
621a43d} and
\href{https://github.com/leanprover-community/mathlib4/tree/0df444a360eaa60ab8c11dca51a86af692955474}{Mathlib
0df444a}. Exact integer-power identities replace fractional-power
approximation. One kernel-mode LeanCert certificate proves the positive
rational gap and remains in the final norm and eigenvalue contradiction.
See the
\href{https://github.com/ajt60gaibb/OpenProblemsInNLA/blob/main/matrix-inequalities-and-norms/MI-23/lean/formalization.yaml}{manifest},
\href{https://github.com/ajt60gaibb/OpenProblemsInNLA/blob/main/matrix-inequalities-and-norms/MI-23/lean/lake-manifest.json}{dependency
pins} and
\href{https://github.com/ajt60gaibb/OpenProblemsInNLA/blob/main/matrix-inequalities-and-norms/MI-23/lean/NUMERICAL_TARGETS.md}{numerical
targets}. From the verified revision on a documented
\href{https://github.com/ajt60gaibb/OpenProblemsInNLA/blob/main/tools/lean/HARNESS.md}{non-root
Linux host}, reproduce with:

\begin{verbatim}
tools/lean/bootstrap.sh /absolute/path/to/nla-lean-tools
tools/lean/selftest.sh /absolute/path/to/nla-lean-tools
tools/lean/verify.sh \
  matrix-inequalities-and-norms/MI-23/lean \
  /absolute/path/to/nla-lean-tools
\end{verbatim}

This settles the corrected eigenvalue conjecture stated below. The
earlier singular-value conjecture is a different target. The original
problem and historical informal proof remain intact.

\subsection{Resolution --- 2026-09-11}\label{resolution-2026-09-11}

\textbf{Negative result by Matthew J. Colbrook}, Department of Applied
Mathematics and Theoretical Physics, University of Cambridge.
\textbf{Independent proof review: PASS.}

Rational positive definite \(3\times3\) matrices with \(r=s=1\), \(p=2\)
and \(t=1/8\) violate the corrected eigenvalue log-majorization
conjecture. An exact integer-power construction and rational norm
separation establish failure of the first ordered eigenvalue inequality.

The exact target is resolved. The original statement and source evidence
are retained below; its former difficulty rating is historical.

\textbf{Primary manuscript:}
\href{https://github.com/ajt60gaibb/OpenProblemsInNLA/blob/main/matrix-inequalities-and-norms/MI-23/solution.pdf}{complete
proof PDF},
\href{https://github.com/ajt60gaibb/OpenProblemsInNLA/blob/main/matrix-inequalities-and-norms/MI-23/solution.tex}{standalone
TeX}, Theorem 1.1 and its proof;
\href{https://github.com/ajt60gaibb/OpenProblemsInNLA/blob/main/matrix-inequalities-and-norms/MI-23/solution.md}{authorship
and scope}. The
\href{https://github.com/ajt60gaibb/OpenProblemsInNLA/blob/main/references/colbrook-matrix-2026-09-11/verification/reviews/MI-23-review.md}{independent
review} checks the full original argument and records its hash. The
draft was AI-assisted; the 2026-09-11 review was independent agent
verification, without external human peer review or formal
certification. The later Lean verification above covers the complete
corrected target.
\href{https://github.com/ajt60gaibb/OpenProblemsInNLA/blob/main/references/colbrook-matrix-2026-09-11/README.md}{Submission
record}.

\subsection{Problem statement}\label{problem-statement}

For positive definite \(A,B\in\mathbb C^{n\times n}\), real \(r\), and
\(t\in[0,1]\), define

\[G_{r,t}(A,B)=A^{r/2}(A^{-1/2}BA^{-1/2})^tA^{r/2}.\]

Is it true, for every \(n\ge1\), every such \(A,B\), every \(p\ge1\),
\(t\in[0,1]\), and either \(r,s\ge1\) or \(r,s\le0\), that

\[\lambda\bigl(G_{r,t}(A,B)^pG_{s,1-t}(A,B)^p\bigr)
\prec_{\log}\lambda\bigl(A^{p(r+s-1)}B^p\bigr)?\]

The eigenvalues of these products are positive real numbers, ordered
decreasingly: each product is similar to a positive definite matrix. For
positive decreasing vectors \(x,y\in\mathbb R^n\), \(x\prec_{\log}y\)
means \(\prod_{j=1}^kx_j\le\prod_{j=1}^ky_j\) for \(1\le k< n\), with
equality at \(k=n\).

This compares nonlinear matrix-mean products with products of powers of
the input matrices. It would give simultaneous multiplicative control
for every partial product of ordered eigenvalues and strengthen norm
estimates used with positive definite matrix functions.

\newpage
\subsection{References}\label{references}

\begin{enumerate}
\def\labelenumi{\arabic{enumi}.}
\tightlist
\item
  M. M. Ghabries, H. Abbas, B. Mourad and A. Assi, \emph{New
  log-majorization results concerning eigenvalues and singular values
  and a complement of a norm inequality}, preprint (2021), Conjecture
  2.1, p.~6. \href{https://arxiv.org/pdf/2105.13356}{PDF}.
\item
  M. M. Ghabries, \emph{Contributions to Matrix Inequalities and Some
  Applications}, PhD thesis, University of Angers and Lebanese
  University (2022), final Open Problems, Problem 6, pp.~109--110.
  \href{https://www.researchgate.net/publication/361793582_Contributions_to_Matrix_Inequalities_and_Some_Applications}{Thesis
  uploaded by its author}.
\item
  M. M. Ghabries, \emph{A log-majorization inequality for normal
  matrices with applications to determinantal inequalities and geometric
  means} (2026), §4. \href{https://arxiv.org/html/2607.21163}{Full
  text}.
\end{enumerate}

Status check (2026-09-10): Theorem 2.1 and equation (13) of the first
source prove a substantive included region: \(r,s\ge1\), \(1\le p\le2\),
and \(\frac{r(p-1)}{(r+s)p}\le t\le\frac{rp+s}{(r+s)p}\). This includes
\(p=1\) for all \(t\). The full parameter range is the remaining
question. The first source refutes its earlier unrestricted
singular-value Conjecture 1.2; the statement here is the explicitly
proposed replacement. The thesis retains it. The July 2026 paper
concerns eigenvalue products of ordinary weighted means and does not
claim this full \(r,s,p\) statement. Current arXiv versions and searches
for the authors, ``Conjecture 2.1 generalized geometric mean'' and
2025/2026 yielded no full resolution. This is a bounded check, not
certification of openness.
\end{document}
